% Nguồn SGK Toán 9 Tập hai — Luyện tập chung sau Bài 21, trang 28--29:
% https://cdn3.olm.vn/upload/thuvienso/4491669493-page-28-1771844543945.png
% https://cdn3.olm.vn/upload/thuvienso/4491669493-page-29-1771844544269.png

\section*{Luyện tập chung}

\begin{lesson}
    \textbf{Ví dụ 1}

    Cho phương trình bậc hai: $x^2-7x+5=0$.
    \begin{enumerate}[label=\alph*)]
        \item Chứng minh phương trình có hai nghiệm phân biệt $x_1,x_2$.
        \item Tính $A=x_1^2+x_2^2$.
    \end{enumerate}

    \textbf{Giải}

    a) Ta có: $\Delta=7^2-4\cdot 5=29>0$ nên phương trình có hai nghiệm phân biệt $x_1,x_2$.

    b) Theo định lí Viète, ta có: $x_1+x_2=7$, $x_1x_2=5$.

    Do đó
    \[
        A=x_1^2+x_2^2=(x_1+x_2)^2-2x_1x_2=7^2-2\cdot 5=39.
    \]

    \textbf{Ví dụ 2}

    Có hai miếng kim loại: miếng thứ nhất nặng 600 g; miếng thứ hai nặng 640 g. Thể tích của miếng thứ nhất lớn hơn thể tích của miếng thứ hai là 120 cm$^3$, nhưng khối lượng riêng của miếng thứ nhất nhỏ hơn khối lượng riêng của miếng thứ hai là 5 g/cm$^3$. Tìm khối lượng riêng của mỗi miếng kim loại.

    \textbf{Giải}

    Gọi $x$ (g/cm$^3$) là khối lượng riêng của miếng kim loại thứ nhất. Điều kiện: $x>0$.

    Khối lượng riêng của miếng kim loại thứ hai là $x+5$ (g/cm$^3$).

    Thể tích của miếng kim loại thứ nhất là $\dfrac{600}{x}$ (cm$^3$).

    Thể tích của miếng kim loại thứ hai là $\dfrac{640}{x+5}$ (cm$^3$).

    Theo đề bài, ta có phương trình:
    \[
        \dfrac{600}{x}-\dfrac{640}{x+5}=120.
    \]

    Nhân cả hai vế của phương trình với $x(x+5)$ để khử mẫu, ta được:
    \[
        600(x+5)-640x=120x(x+5),
    \]
    hay $3x^2+16x-75=0$.

    Ta có: $\Delta'=8^2-3\cdot(-75)=289$; $\sqrt{\Delta'}=17$.

    Suy ra phương trình có hai nghiệm phân biệt:
    \[
        x_1=\dfrac{-8-17}{3}=\dfrac{-25}{3}\ \text{(loại)};
        \qquad
        x_2=\dfrac{-8+17}{3}=3\ \text{(thỏa mãn điều kiện)}.
    \]

    Vậy khối lượng riêng của miếng kim loại thứ nhất là 3 g/cm$^3$ và khối lượng riêng của miếng kim loại thứ hai là 8 g/cm$^3$.
\end{lesson}

\begin{exercises}
    \subsection{Bài tập}

    \begin{questions}[resume=SGK]
        \item Tính nhẩm nghiệm của các phương trình sau:
        \begin{questions}
            \item $\sqrt{2}x^2-(\sqrt{2}+1)x+1=0$;
            \item $2x^2+(\sqrt{3}-1)x-3+\sqrt{3}=0$.
        \end{questions}

        \answerline
        \answerline

        \item Gọi $x_1,x_2$ là hai nghiệm của phương trình bậc hai $x^2-5x+3=0$. Không giải phương trình, hãy tính:
        \begin{questions}
            \item $x_1^2+x_2^2$;
            \item $(x_1-x_2)^2$.
        \end{questions}

        \answerline
        \answerline
        \answerline

        \item Tìm hai số $u$ và $v$, biết:
        \begin{questions}
            \item $u+v=15$, $uv=56$;
            \item $u^2+v^2=125$, $uv=22$.
        \end{questions}

        \answerline
        \answerline
        \answerline

        \item Một chiếc hộp có dạng hình hộp chữ nhật, không có nắp, có đáy là hình vuông, tổng diện tích xung quanh và diện tích đáy là 800 cm$^2$. Chiều cao của hộp là 10 cm. Tính độ dài cạnh đáy của chiếc hộp (làm tròn kết quả đến hàng phần mười của cm).

        \answerline
        \answerline
        \answerline

        \item Nhu cầu của khách hàng đối với một loại áo phông tại một cửa hàng được cho bởi phương trình $p=100-0{,}02x$, trong đó $p$ là giá tiền của mỗi chiếc áo (nghìn đồng) và $x$ là số lượng áo phông bán được. Doanh thu $R$ (nghìn đồng) khi bán được $x$ chiếc áo phông là:
        \[
            R=xp=x(100-0{,}02x).
        \]
        Hỏi cần phải bán được bao nhiêu chiếc áo phông để doanh thu đạt 120 triệu đồng?

        \answerline
        \answerline
        \answerline
    \end{questions}
\end{exercises}
